\documentclass[10pt]{extarticle}
\usepackage[a4paper,margin=1.5cm]{geometry}
\usepackage[T1]{fontenc}
\usepackage{amsmath,amssymb}
\usepackage{graphicx}
\usepackage[export]{adjustbox}
\usepackage{multicol}
\usepackage{enumitem}
\usepackage{booktabs,array,tabularx}
\setlength{\columnsep}{0.7cm}
\setlength{\parindent}{0pt}
\setlength{\parskip}{3pt}
\pagestyle{empty}
\makeatletter
% figure registry (figdefs.tex is generated by the extract step)
\newcommand{\qfigdef}[4]{\@namedef{qf@file@#1}{#2}\@namedef{qf@w@#1}{#3}\@namedef{qf@wide@#1}{#4}}
\newcommand{\qf@check}[1]{\@ifundefined{qf@file@#1}{\PackageError{score2}{Unknown figure id #1}{}}{}}
\newcommand{\qf@img}[4][]{\includegraphics[width=#2\dimexpr\@nameuse{qf@w@#4}pt\relax,max width=#3,max height=0.45\textheight,#1]{\@nameuse{qf@file@#4}}}
% Full-width questions: finalize wraps every question that contains a wide figure in \qfullbegin...\qfullend,
% so the WHOLE question (stem, figures, options) is set across the page in reading order, then the two columns
% resume. The question's head - its number, stem and first figure (\qfullsplit marks the end of the head) -
% is kept on one page, so a number is never left at a page foot with its figure overleaf; the rest may break
% across pages normally, so no big gaps are left. Without a \qfullsplit, the whole question is kept together
% if it is under 45% of a page. Never write these yourself.
\newif\ifqf@full
\newif\ifqf@boxed
\newsavebox{\qf@box}
\newcommand{\qf@open}{\global\qf@fulltrue\par\global\qf@boxedtrue
  \setbox\qf@box\vbox\bgroup\hsize\textwidth\linewidth\textwidth}
\newcommand{\qf@emit}[1]{\par\egroup\global\qf@boxedfalse
  \ifdim\dimexpr\ht\qf@box+\dp\qf@box\relax<#1\textheight\noindent\usebox{\qf@box}\par\else\unvbox\qf@box\fi}
\newcommand{\qfullsplit}{\ifqf@boxed\qf@emit{0.85}\fi}
\newcommand{\qf@close}{\ifqf@boxed\qf@emit{0.45}\fi\global\qf@fullfalse}
\newcommand{\qfullbegin}{\end{multicols}\qf@open}
\newcommand{\qfullend}{\qf@close\begin{multicols}{2}}
\newcommand{\qfullstart}{\qf@open}
\newcommand{\qfulldone}{\qf@close}
% \qfig[scale]{id}: block figure, centred, directly where it is written (size set by the tool).
% Fallback if a wide figure is not inside a full-width question: break out of the columns in place - no floats.
\newcommand{\qfig}[2][1]{\qf@check{#2}%
  \ifqf@full\qf@blk{#1}{#2}\else
  \ifnum\@nameuse{qf@wide@#2}=1 \expandafter\@firstoftwo\else\expandafter\@secondoftwo\fi
  {\end{multicols}\begin{center}\qf@img{#1}{\textwidth}{#2}\end{center}\begin{multicols}{2}}%
  {\qf@blk{#1}{#2}}\fi}
\newcommand{\qf@blk}[2]{\par\smallskip{\centering\qf@img{#1}{\linewidth}{#2}\par}\smallskip}
% \qfigi[scale]{id}: inline figure (inside option lists / tables / minipages); never breaks columns.
\newcommand{\qfigi}[2][1]{\qf@check{#2}\qf@img[valign=c]{#1}{\linewidth}{#2}}
\makeatother
% structure macros - lint reads these to check completeness and figure placement
\newcommand{\papersec}[1]{\par\bigskip{\centering\Large\bfseries PAPER-#1\par}\medskip}
\newcommand{\partsec}[2]{\par\medskip\noindent{\large\bfseries PART-#1: #2}\par\smallskip}
\newcommand{\sect}[1]{\par\smallskip\noindent\textbf{#1}\par}
\newcommand{\qn}[1]{\par\smallskip\noindent\textbf{#1.}~}
\newcommand{\paraq}[2]{\par\smallskip\noindent\textbf{Paragraph for Questions #1 and #2}\par}
\newenvironment{qoptions}{\begin{itemize}[leftmargin=1.8em,itemsep=1pt,topsep=2pt,label={}]}{\end{itemize}}
\newcommand{\opt}[2]{\item[(#1)] #2}
% figure registry (generated by score2_tools.py finalize)
\qfigdef{p1_phy_s1_q01_1}{Score Advanced - Test 24 (17 Apr 2024)-img/p1_phy_s1_q01_1.png}{76.8}{0}
\qfigdef{p1_phy_s1_q02_1}{Score Advanced - Test 24 (17 Apr 2024)-img/p1_phy_s1_q02_1.png}{88.9}{0}
\qfigdef{p1_phy_s1_q03_1}{Score Advanced - Test 24 (17 Apr 2024)-img/p1_phy_s1_q03_1.png}{90.6}{0}
\qfigdef{p1_phy_s1_q05_1}{Score Advanced - Test 24 (17 Apr 2024)-img/p1_phy_s1_q05_1.png}{123.5}{0}
\qfigdef{p1_phy_s2_pre_1}{Score Advanced - Test 24 (17 Apr 2024)-img/p1_phy_s2_pre_1.png}{163.2}{0}
\qfigdef{p1_phy_s2_pre_2}{Score Advanced - Test 24 (17 Apr 2024)-img/p1_phy_s2_pre_2.png}{170.8}{0}
\qfigdef{p1_phy_s2_pre_3}{Score Advanced - Test 24 (17 Apr 2024)-img/p1_phy_s2_pre_3.png}{168.0}{0}
\qfigdef{p1_phy_s2_pre_4}{Score Advanced - Test 24 (17 Apr 2024)-img/p1_phy_s2_pre_4.png}{167.5}{0}
\qfigdef{p1_phy_s2_q09_1}{Score Advanced - Test 24 (17 Apr 2024)-img/p1_phy_s2_q09_1.png}{135.6}{0}
\qfigdef{p1_phy_s3_q01_1}{Score Advanced - Test 24 (17 Apr 2024)-img/p1_phy_s3_q01_1.png}{81.8}{0}
\qfigdef{p1_phy_s3_q03_1}{Score Advanced - Test 24 (17 Apr 2024)-img/p1_phy_s3_q03_1.png}{106.6}{0}
\qfigdef{p1_phy_s3_q04_1}{Score Advanced - Test 24 (17 Apr 2024)-img/p1_phy_s3_q04_1.png}{111.9}{0}
\qfigdef{p1_phy_s3_q05_1}{Score Advanced - Test 24 (17 Apr 2024)-img/p1_phy_s3_q05_1.png}{147.4}{0}
\qfigdef{p2_phy_s1_q01_1}{Score Advanced - Test 24 (17 Apr 2024)-img/p2_phy_s1_q01_1.png}{115.1}{0}
\qfigdef{p2_phy_s1_q03_1}{Score Advanced - Test 24 (17 Apr 2024)-img/p2_phy_s1_q03_1.png}{98.7}{0}
\qfigdef{p2_phy_s2_q06_1}{Score Advanced - Test 24 (17 Apr 2024)-img/p2_phy_s2_q06_1.png}{93.3}{0}
\qfigdef{p2_phy_s2_q09_1}{Score Advanced - Test 24 (17 Apr 2024)-img/p2_phy_s2_q09_1.png}{77.4}{0}
\qfigdef{p2_phy_s2_q09_A}{Score Advanced - Test 24 (17 Apr 2024)-img/p2_phy_s2_q09_A.png}{70.3}{0}
\qfigdef{p2_phy_s2_q09_B}{Score Advanced - Test 24 (17 Apr 2024)-img/p2_phy_s2_q09_B.png}{68.6}{0}
\qfigdef{p2_phy_s2_q09_C}{Score Advanced - Test 24 (17 Apr 2024)-img/p2_phy_s2_q09_C.png}{70.8}{0}
\qfigdef{p2_phy_s2_q09_D}{Score Advanced - Test 24 (17 Apr 2024)-img/p2_phy_s2_q09_D.png}{71.3}{0}
\qfigdef{p2_phy_s2_q10_1}{Score Advanced - Test 24 (17 Apr 2024)-img/p2_phy_s2_q10_1.png}{59.3}{0}
\qfigdef{p2_phy_s3_q02_1}{Score Advanced - Test 24 (17 Apr 2024)-img/p2_phy_s3_q02_1.png}{101.1}{0}
\qfigdef{p2_phy_s3_q03_1}{Score Advanced - Test 24 (17 Apr 2024)-img/p2_phy_s3_q03_1.png}{111.7}{0}
\qfigdef{p2_phy_s3_q04_1}{Score Advanced - Test 24 (17 Apr 2024)-img/p2_phy_s3_q04_1.png}{157.0}{0}
\qfigdef{p2_phy_s3_q07_1}{Score Advanced - Test 24 (17 Apr 2024)-img/p2_phy_s3_q07_1.png}{106.1}{0}
\qfigdef{p2_chem_s1_q02_1}{Score Advanced - Test 24 (17 Apr 2024)-img/p2_chem_s1_q02_1.png}{189.3}{0}
\qfigdef{p2_chem_s1_q02_A}{Score Advanced - Test 24 (17 Apr 2024)-img/p2_chem_s1_q02_A.png}{133.4}{0}
\qfigdef{p2_chem_s1_q02_B}{Score Advanced - Test 24 (17 Apr 2024)-img/p2_chem_s1_q02_B.png}{132.0}{0}
\qfigdef{p2_chem_s1_q02_C}{Score Advanced - Test 24 (17 Apr 2024)-img/p2_chem_s1_q02_C.png}{131.1}{0}
\qfigdef{p2_chem_s1_q02_D}{Score Advanced - Test 24 (17 Apr 2024)-img/p2_chem_s1_q02_D.png}{133.5}{0}
\qfigdef{p2_chem_s2_q07_A}{Score Advanced - Test 24 (17 Apr 2024)-img/p2_chem_s2_q07_A.png}{51.8}{0}
\qfigdef{p2_chem_s2_q07_B}{Score Advanced - Test 24 (17 Apr 2024)-img/p2_chem_s2_q07_B.png}{51.1}{0}
\qfigdef{p2_chem_s2_q07_C}{Score Advanced - Test 24 (17 Apr 2024)-img/p2_chem_s2_q07_C.png}{32.3}{0}
\qfigdef{p2_chem_s2_q07_D}{Score Advanced - Test 24 (17 Apr 2024)-img/p2_chem_s2_q07_D.png}{50.4}{0}
\qfigdef{p2_chem_s2_q09_C}{Score Advanced - Test 24 (17 Apr 2024)-img/p2_chem_s2_q09_C.png}{125.2}{0}
\qfigdef{p2_chem_s2_q09_D}{Score Advanced - Test 24 (17 Apr 2024)-img/p2_chem_s2_q09_D.png}{119.0}{0}
\qfigdef{p2_math_s2_q06_1}{Score Advanced - Test 24 (17 Apr 2024)-img/p2_math_s2_q06_1.png}{54.4}{0}
\qfigdef{p2_math_s3_q03_1}{Score Advanced - Test 24 (17 Apr 2024)-img/p2_math_s3_q03_1.png}{182.4}{0}
\raggedcolumns
\makeatletter\@beginparpenalty=10000\makeatother
\newcolumntype{L}[1]{>{\raggedright\arraybackslash}p{#1}}
\begin{document}
\begin{center}{\Large\bfseries Score Advanced --- Test 24 (17 Apr 2024)}\end{center}
\begin{multicols}{2}
\raggedcolumns
\papersec{1}
\partsec{1}{PHYSICS}
\sect{SECTION-I (i)}
\qn{1} Two vertical conducting cylinders A and B of different cross sections are connected by a thin tube as shown. An ideal gas is trapped in the cylinders by two airtight pistons of masses 1.0 kg and 2.0 kg respectively. Initially the pistons are at the same height 40 cm above the base of the cylinders. If an additional mass of 10 g were gently placed on the piston in the cylinder A, which of the following statements correctly describe the new steady state ultimately reached.\par
\qfig{p1_phy_s1_q01_1}
\begin{qoptions}
\opt{A}{Pressure of the gas has been increased.}
\opt{B}{Pressure of the gas remains unchanged.}
\opt{C}{Height difference between the pistons becomes 30 cm.}
\opt{D}{Height difference between the pistons becomes 60 cm.}
\end{qoptions}
\qn{2} Five equal resistances, a capacitor and two batteries are connected in a circuit as shown in the figure. The capacitor is initially uncharged and each battery has negligible internal resistance. It is given that R = 1$\Omega$, C = 1$\mu$F and V = 8 volt. On the basis of above information answer the following questions, if switch `S' is closed at t = 0. (Assume $e^{-1}=0.37$, if required)\par
\qfig{p1_phy_s1_q02_1}
\begin{qoptions}
\opt{A}{Current in the branch of capacitor at t = 5/8 $\mu$s is (approximately) 6 A}
\opt{B}{Current in the branch AB at t = 5/8 $\mu$s is (approximately) 1.7 A}
\opt{C}{Potential difference between points A and C at t = 5/8 $\mu$s is (approximately) 2.5 V}
\opt{D}{Potential difference between points A and C at t = 5/8 $\mu$s is (approximately) 6 V}
\end{qoptions}
\qn{3} A point mass m is attached to an ideal string, which in turns is wrapped on a vertical pole. Point mass is kept on a smooth horizontal surface and string is horizontal. A horizontal velocity is imparted to the point mass perpendicular to length of string.\par
\qfig{p1_phy_s1_q03_1}
\begin{qoptions}
\opt{A}{Speed of mass will become twice of initial value if length of string between the point mass and the pole becomes half of the original value.}
\opt{B}{Kinetic energy of mass increases.}
\opt{C}{Work done by tension on the point mass is zero.}
\opt{D}{Energy of the mass remains constant.}
\end{qoptions}
\qn{4} In the adjacent figure, a uniform disc of mass 2m and radius $\ell/2$ is lying at rest on a smooth horizontal surface. A particle `A' of mass m is connected to a light string of length $\ell$, whose other end is attached to the circumference of the disc. Initially string is just taut and tangential to the disc, particle A is at rest. In the same horizontal plane another particle B of same mass m moving with velocity $v_0$ perpendicular to string collides elastically with A. Just after impact which of the following statements will be true
\begin{qoptions}
\opt{A}{Tension in the string is $\dfrac{2mv_0^2}{5\ell}$}
\opt{B}{Acceleration of the centre of the disc is $\dfrac{v_0^2}{5\ell}$}
\opt{C}{Initial radius of curvature of `A' is $\dfrac{5\ell}{2}$}
\opt{D}{Acceleration of the centre of the disc is $\dfrac{2v_0^2}{5\ell}$}
\end{qoptions}
\qn{5} Three masses are joined using pin joints (can move freely). Lengths of both rods are identical and the system starts from rest. Accelerations of the balls are shown in the figure. The angle between the rods at this moment is $90^\circ$. The horizontal surface is smooth. Choose the correct option(s).\par
\qfig{p1_phy_s1_q05_1}
\begin{qoptions}
\opt{A}{$2a_2=a_1+a_x$}
\opt{B}{$2a_y=a_1+a_2$}
\opt{C}{$2a_x=a_1-a_2$}
\opt{D}{$3a_1=5a_2$}
\end{qoptions}
\qn{6} A particle of mass m is subjected to two forces
\[\vec F_1=f(r)\dfrac{\vec r}{|\vec r|}\quad\text{and}\quad \vec F_2=-\alpha\vec v\]
where f(r) is a real valued function depending on the distance of the particle from the origin, $\vec r$ is the position vector of the particle and $\alpha$ is a positive constant. If at t = 0, $\vec v$ vector is in the x-y plane and the angular momentum of the particle about the origin is $L_0$. Choose the correct option(s).
\begin{qoptions}
\opt{A}{Plane of motion of the particle does not change.}
\opt{B}{Angular momentum of the particle about the origin is given by $L_0e^{-\frac{\alpha t}{m}}$.}
\opt{C}{Angular momentum of the particle about the origin remains constant.}
\opt{D}{The value of angular momentum about the origin will depend on f(r).}
\end{qoptions}
\sect{SECTION-I (ii)}
\qn{7} List-I (Circuit)\par
(I)\par\qfig{p1_phy_s2_pre_1}
(II)\par\qfig{p1_phy_s2_pre_2}
(III)\par\qfig{p1_phy_s2_pre_3}
(IV)\par\qfig{p1_phy_s2_pre_4}
List-II ($C_{eq}$): (P) $\dfrac{3C}{4}$; (Q) $\dfrac{5C}{12}$; (R) $\dfrac{2C}{3}$; (S) $\dfrac{3C}{2}$
\begin{qoptions}
\opt{A}{I $\rightarrow$ S;II $\rightarrow$ Q;III $\rightarrow$ R;IV $\rightarrow$ P}
\opt{B}{I $\rightarrow$ S;II $\rightarrow$ P;III $\rightarrow$ Q;IV $\rightarrow$ R}
\opt{C}{I $\rightarrow$ S;II $\rightarrow$ P;III $\rightarrow$ R;IV $\rightarrow$ Q}
\opt{D}{I $\rightarrow$ S;II $\rightarrow$ Q;III $\rightarrow$ P;IV $\rightarrow$ R}
\end{qoptions}
\qn{8} One mole of an ideal gas whose adiabatic exponent equals $\gamma$ undergoes a process $P=P_0+\alpha/v$, where $P_0$ and $\alpha$ are positive constants. Its volume increases from $v_1$ to $v_2$.\par
{\footnotesize\setlength{\tabcolsep}{1.5pt}\noindent\begin{tabular}{|c|L{3.0cm}|c|L{3.2cm}|} \hline & \multicolumn{1}{c|}{List-I} & & \multicolumn{1}{c|}{List-II}\\ \hline
(I) & Heat capacity of gas as a function of its volume & (P) & $\dfrac{P_0(V_2-V_1)}{\gamma-1}$\\ \hline
(II) & Internal energy increment of a gas & (Q) & $\dfrac{P_0(V_2-V_1)}{\gamma-1}+\alpha\ln\dfrac{V_2}{V_1}$\\ \hline
(III) & The work performed by it & (R) & $P_0(V_2-V_1)+\alpha\ln\dfrac{V_2}{V_1}$\\ \hline
(IV) & The amount of heat transferred & (S) & $\dfrac{\gamma R}{\gamma-1}+\dfrac{\alpha R}{p_0v}$\\ \hline
\end{tabular}}\par
\begin{qoptions}
\opt{A}{I $\rightarrow$ S;II $\rightarrow$ P;III $\rightarrow$ R;IV $\rightarrow$ Q}
\opt{B}{I $\rightarrow$ P;II $\rightarrow$ Q;III $\rightarrow$ R;IV $\rightarrow$ S}
\opt{C}{I $\rightarrow$ P;II $\rightarrow$ S;III $\rightarrow$ R;IV $\rightarrow$ Q}
\opt{D}{I $\rightarrow$ S;II $\rightarrow$ P;III $\rightarrow$ Q;IV $\rightarrow$ S}
\end{qoptions}
\qn{9} Consider a situation in which a projectile is fired at 40 m/s from the top of a tower in absence of air resistance, t = 0.\par
\qfig{p1_phy_s2_q09_1}
{\footnotesize\setlength{\tabcolsep}{1.5pt}\noindent\begin{tabular}{|c|L{3.9cm}|c|L{1.9cm}|} \hline & \multicolumn{1}{c|}{List-I} & & \multicolumn{1}{c|}{List-II}\\ \hline
(I) & Time at which particle moves at right angles with the initial direction & (P) & 2.4 s\\ \hline
(II) & time at which kinetic energy is least & (Q) & 6 s\\ \hline
(III) & Time at which it strikes the ground & (R) & 20/3 s\\ \hline
(IV) & Time when $\vec v$ make $45^\circ$ with horizontal & (S) & 5.6 s\\ \hline
 & & (T) & 7.2 s\\ \hline
\end{tabular}}\par
\begin{qoptions}
\opt{A}{I $\rightarrow$ R;II $\rightarrow$ P;III $\rightarrow$ S;IV $\rightarrow$ Q}
\opt{B}{I $\rightarrow$ P;II $\rightarrow$ R;III $\rightarrow$ Q;IV $\rightarrow$ S}
\opt{C}{I $\rightarrow$ P;II $\rightarrow$ Q;III $\rightarrow$ R;IV $\rightarrow$ S}
\opt{D}{I $\rightarrow$ R;II $\rightarrow$ P;III $\rightarrow$ Q;IV $\rightarrow$ S}
\end{qoptions}
\qn{10} A neutron moving with kinetic energy K collides with a hydrogen atom at rest in ground state. Assume mass of neutron to be equal to that of hydrogen atom. Match the kinetic energy of neutron with the possible phenomena mentioned in List-II. Energy of hydrogen in ground state is $-13.6$ eV. Match List-I with List-II :\par
{\footnotesize\setlength{\tabcolsep}{1.5pt}\noindent\begin{tabular}{|c|L{2.4cm}|c|L{3.0cm}|} \hline & \multicolumn{1}{c|}{\parbox{2.4cm}{\centering List-I (Kinetic energy of colliding neutron)}} & & \multicolumn{1}{c|}{List-II}\\ \hline
(I) & K = 10.2 eV & (P) & Elastic collision may take place\\ \hline
(II) & K = 13.6 eV & (Q) & Inelastic collision may take place\\ \hline
(III) & K = 20.4 eV & (R) & Perfectly inelastic collision may take place\\ \hline
(IV) & K = 22.0 eV & (S) & Hydrogen atom may get ionised\\ \hline
\end{tabular}}\par
\begin{qoptions}
\opt{A}{I $\rightarrow$ P;II $\rightarrow$ Q;III $\rightarrow$ R;IV $\rightarrow$ S}
\opt{B}{I $\rightarrow$ P;II $\rightarrow$ P;III $\rightarrow$ P,R;IV $\rightarrow$ P,Q}
\opt{C}{I $\rightarrow$ P;II $\rightarrow$ P;III $\rightarrow$ P;IV $\rightarrow$ Q}
\opt{D}{I $\rightarrow$ P;II $\rightarrow$ P;III $\rightarrow$ R;IV $\rightarrow$ Q}
\end{qoptions}
\sect{SECTION-II}
\qn{1} Two identical springs each of stiffness k are welded at a point P. If a particle of mass m is welded at P. The period of oscillation of the particle in the direction of x is $\dfrac{\pi}{\sin\theta}\sqrt{\dfrac{Ym}{k}}$. Find Y.\par
\qfig{p1_phy_s3_q01_1}
\qn{2} A body executing SHM with time period T and amplitude A. When it is at $\dfrac A2$ from equilibrium position moving toward right, it receives an impulse which doubles its velocity without changing direction. Time taken by the block to reach the right extreme position is $t=T\cos^{-1}\left(\dfrac{K}{\sqrt{13}}\right)$. Find K.
\qn{3} Figure shows a meter bridge wire AC has a uniform cross section. The length of wire AC is 100 cm. X is a standard resistor of 4 $\Omega$ and Y is a coil. When Y is immersed in melting ice the null point is at 40 cm from point A. When the coil Y is heated to $100^\circ$C, a 12 $\Omega$ resistor has to be connected in parallel with Y in order to keep the bridge balanced at the same point.\par
\qfig{p1_phy_s3_q03_1}
\qn{4} A circular disc of a diameter d is slowly rotated in a liquid of large viscosity $\mu$ at a small distance h from a fixed surface. An expression for torque T necessary to maintain an angular velocity $\omega$ is $\dfrac{\mu\pi d^4\omega}{Xh}$. Find the value of X.\par
\qfig{p1_phy_s3_q04_1}
\qn{5} A set of parallel rays incident on the bi-convex lens. If after two refractions and one reflection final image formed at 15 cm left of lens, the value of $\dfrac X9$ is\par
\qfig{p1_phy_s3_q05_1}
\qn{6} Three SHM in the same direction in having same amplitude and same period are superposed. If each differ in phase from the next by $45^\circ$. Then, find the energy associated with resulting motion taking energy associated with each of given SHM as 1 J.
\qn{7} Two rods of different materials having coefficient of thermal expansion $\alpha_1,\alpha_2$ and Young's modulus $Y_1,Y_2$ are fixed between two massive walls. The rods are heated such that they undergo the same increase in temperature. If $\alpha_1:\alpha_2=2:3$, then thermal stresses in rods are equal, then the ratio $Y_1:Y_2$ is \_\_\_\_\_.
\qn{8} A student determines a dimensionless quantity,
\[B=\dfrac{e^n}{Z\varepsilon_0hc}.\]
Where, e = electric charge,\par
$\varepsilon_0$ = permittivity of vacuum,\par
h = Plank's constant\par
and c = speed of light.\par
Find the value of n.
\partsec{2}{CHEMISTRY}
\sect{SECTION-I (i)}
\qn{1} An isotone of $^{76}_{32}$Ge is
\begin{qoptions}
\opt{A}{$^{77}_{32}$Ge}
\opt{B}{$^{77}_{33}$As}
\opt{C}{$^{77}_{34}$Se}
\opt{D}{$^{78}_{34}$Se}
\end{qoptions}
\qn{2} For a compound AB crystallises in a rock salt structure with edge length a of the unit cell,
\begin{qoptions}
\opt{A}{the minimum distance between A and A is $a/\sqrt2$}
\opt{B}{the minimum distance between B and B is $a/\sqrt2$}
\opt{C}{the minimum distance between A and B is $a/2$}
\opt{D}{the density of the crystal is given by $\rho=\dfrac{4}{a^3}\dfrac{M_{AB}}{N_A}$}
\end{qoptions}
\qn{3} Which of the following statements are not correct
\begin{qoptions}
\opt{A}{For an ideal binary liquid solution, the plot of $1/x_A$ versus $1/y_A$ (where $x_A$ and $y_A$ are the mole fractions of A in liquid and vapour phases, respectively) is linear.}
\opt{B}{According to Raoult's law, the relative lowering of vapour pressure is equal to the mole fraction of solvent in the solution.}
\opt{C}{The crysoscopic constant $K_f$ is property of solvent.}
\opt{D}{The value of van't Hoff factor for CH$_3$COOH in benzene will be 2}
\end{qoptions}
\qn{4} Which of the following statements are correct?
\begin{qoptions}
\opt{A}{Linear overlap of two atomic p-orbitals leads to a sigma bond}
\opt{B}{The bond angle H $-$ N $-$ H in NH$_3$ is greater than the bond angle H $-$ As $-$ H in AsH$_3$.}
\opt{C}{Anhydrous HCl is a bad conductor of electricity but aqueous HCl is a good conductor}
\opt{D}{o-nitrophenol is steam volatile whereas p-nitrophenol is not.}
\end{qoptions}
\qn{5} Which of the following reactions give an alkyne?
\begin{qoptions}
\opt{A}{sodium fumarate $\xrightarrow[\text{electrolysis}]{\text{Kolbe's}}$}
\opt{B}{CH$_3$CBr$_2$CHBr$_2$ $\xrightarrow{\text{Zn/alcohol}/\Delta}$}
\opt{C}{CH$_3$CH$_2$CHBr$_2$ $\xrightarrow{\text{NaNH}_2/\Delta}$}
\opt{D}{CH$_3$CHBr $-$ CHBr $\xrightarrow{\text{NaNH}_2}$}
\end{qoptions}
\qn{6} Which of the following reactions give n-propyl alcohol as the major product ?
\begin{qoptions}
\opt{A}{CH$_3-$CH$=$CH$_2$ + H$_2$O $\xrightarrow{\text{H}^+}$}
\opt{B}{HCHO $\xrightarrow[\text{(ii) hydrolysis}]{\text{(i) CH}_3\text{CH}_2\text{MgBr}}$}
\opt{C}{CH$_3-$CH$=$CH$_2$ $\xrightarrow{\text{cold(aq)/KMnO}_4}$}
\opt{D}{CH$_2=$CH$-$CHO + H$_2$ $\xrightarrow{\text{Ni}}$}
\end{qoptions}
\sect{SECTION-I (ii)}
\qn{7} Match the entries in List-I for the spectral series of hydrogen atom with those given in List-II. The symbols $n_1$ and $n_2$ are the quantum level in the expression\par
{\footnotesize\setlength{\tabcolsep}{1.5pt}\noindent\begin{tabular}{|c|L{2.4cm}|c|L{3.9cm}|} \hline & \multicolumn{1}{c|}{List-I} & & \multicolumn{1}{c|}{List-II}\\ \hline
(I) & Lyman series & (P) & $n_1=2$ and $n_2=3,4,5,\ldots$\\ \hline
(II) & Balmer series & (Q) & $n_1=4$ and $n_2=5,6,7$\\ \hline
(III) & Paschen series & (R) & $n_1=5$ and $n_2=6,7,\ldots$\\ \hline
(IV) & Pfund series & (S) & $n_1=1$ and $n_2=2,3,4,\ldots$\\ \hline
 & & (T) & $n_1=3$ and $n_2=4,5,6,\ldots$\\ \hline
\end{tabular}}\par
Which of the following is correct match
\begin{qoptions}
\opt{A}{I $\rightarrow$ S}
\opt{B}{II $\rightarrow$ Q}
\opt{C}{III $\rightarrow$ R}
\opt{D}{IV $\rightarrow$ T}
\end{qoptions}
\qn{8} Match the List-I and List- II\par
{\footnotesize\setlength{\tabcolsep}{1.5pt}\noindent\begin{tabular}{|c|L{2.3cm}|c|L{3.9cm}|} \hline & \multicolumn{1}{c|}{List-I} & & \multicolumn{1}{c|}{List-II}\\ \hline
(I) & Electrophilic substitution & (P) & C$_6$H$_5$N$_2$Cl + C$_6$H$_5$OH $\xrightarrow{\text{NaOH}}$\\ \hline
(II) & Nucleophilic substitution & (Q) & Toluene + NBS $\rightarrow$\\ \hline
(III) & Free radical substitution & (R) & Dow's process\\ \hline
(IV) & Coupling reaction & (S) & Phenol $\rightarrow$ picric acid\\ \hline
 & & (T) & C$_6$H$_5$ONa + CH$_3$Br $\rightarrow$\\ \hline
\end{tabular}}\par
Which of the following is correct match
\begin{qoptions}
\opt{A}{I $\rightarrow$ R}
\opt{B}{II $\rightarrow$ P}
\opt{C}{III $\rightarrow$ Q}
\opt{D}{IV $\rightarrow$ T}
\end{qoptions}
\qn{9} Match the List-I and List-II\par
{\footnotesize\setlength{\tabcolsep}{1.5pt}\noindent\begin{tabular}{|c|L{2.6cm}|c|L{2.6cm}|} \hline & \multicolumn{1}{c|}{List-I} & & \multicolumn{1}{c|}{List-II}\\ \hline
(I) & 2p$\pi-$3p$\pi$ & (P) & F$_2$C = O\\ \hline
(II) & 3p$\pi-$3d$\pi$ & (Q) & F$_3$P = O\\ \hline
(III) & 2p$\pi-$3d$\pi$ & (R) & F$_3$P = S\\ \hline
(IV) & 2p$\pi-$2d$\pi$ & (S) & F$_2$S = O\\ \hline
 & & (T) & F$_2$C = S\\ \hline
\end{tabular}}\par
Which of the following is correct match
\begin{qoptions}
\opt{A}{I $\rightarrow$ R}
\opt{B}{II $\rightarrow$ T}
\opt{C}{III $\rightarrow$ Q}
\opt{D}{IV $\rightarrow$ P}
\end{qoptions}
\qn{10} List-II give the values of $\gamma$ ($-$ $C_{p,m}/C_{V,m}$) at high temperature of gases which are shown in List-I. Identify the gas with the corresponding $\gamma$ values.\par
{\footnotesize\setlength{\tabcolsep}{1.5pt}\noindent\begin{tabular}{|c|L{2.6cm}|c|L{1.6cm}|} \hline & \multicolumn{1}{c|}{List-I} & & \multicolumn{1}{c|}{List-II}\\ \hline
(I) & Monatomic & (P) & 1.154\\ \hline
(II) & Diatomic & (Q) & 1.286\\ \hline
(III) & Linear Triatomic & (R) & 1.167\\ \hline
(IV) & Nonlinear Triatomic & (S) & 1.407\\ \hline
 & & (T) & 1.667\\ \hline
\end{tabular}}\par
Which of following is correct match
\begin{qoptions}
\opt{A}{I $\rightarrow$ Q}
\opt{B}{II $\rightarrow$ T}
\opt{C}{III $\rightarrow$ P}
\opt{D}{IV $\rightarrow$ R}
\end{qoptions}
\sect{SECTION-II}
\qn{1} The ionization potential of hydrogen atom is 13.6 eV. The energy in electron volt required to remove an electron in the n = 2 state of the hydrogen atom is
\qn{2} The density of a gas is 1.5 g L$^{-1}$ at STP. Its density at 27 $^\circ$C and 1000 Torr would be
\qn{3} Ammonium chloride crystallizes in a body-centred cubic lattice with edge length of 387 pm. If the size of NH$_4^+$ is 154 pm, the size of Cl$^-$ ion would be
\qn{4} A container of volume `V' litre contains an equilibrium mixture which consists of 2 mol each of PCl$_5$, PCl$_3$ and Cl$_2$. The equilibrium pressure is 3 atm and temperature is `T' kelvin. A certain amount of Cl$_2$ gas is now introduced into the container keeping the temperature and presence constant, until the equilibrium volume becomes `2V' litre. How many moles of Cl$_2$ were added ?
\qn{5} The conductivity of a saturated solution of silver oxalate is $4.5\times10^{-5}$ S cm$^{-1}$. If its $K_{sp}=1.35\times10^{-11}$ M$^3$, the molar conductivity of the saturated solution would be
\qn{6} How many grams of potassium perchlorate (KClO$_4$) must be mixed with 108 g of water in order to obtain a solution, which has a mole fraction of KClO$_4$ equal to $\dfrac17$ ?
\qn{7} Calculate effective nuclear charge on last electron of Cl. (using Slater rule)
\qn{8} The number of isomers possible for complex K$_2$[PtClBr$_2$SCN] :
\partsec{3}{MATHEMATICS}
\sect{SECTION-I(i)}
\qn{1} A point P $(\sqrt3,1)$ moves on the circle $x^2+y^2=4$ and after covering a quarter of the circle leaves it tangentially. The equation of a line along which the point moves after leaving the circle is
\begin{qoptions}
\opt{A}{$y=\sqrt3x+4$}
\opt{B}{$\sqrt3y=x+4$}
\opt{C}{$\sqrt3y=x-4$}
\opt{D}{$y=\sqrt3x-4$}
\end{qoptions}
\qn{2} A chord PP$'$ of a parabola cuts the axis of the parabola at O. The feet of the perpendiculars from P and P$'$ on the axis are M and M$'$ respectively. If V is the vertex then VM, VO, VM$'$ are in
\begin{qoptions}
\opt{A}{AP}
\opt{B}{GP}
\opt{C}{HP}
\opt{D}{none of these}
\end{qoptions}
\qn{3} $\alpha,\beta$ and $\gamma$ are three angles given by $\alpha=2\tan^{-1}\left(\sqrt2-1\right)$, $\beta=3\sin^{-1}\dfrac1{\sqrt2}+\sin^{-1}\left(-\dfrac12\right)$ and $\gamma=\cos^{-1}\dfrac13$. Then
\begin{qoptions}
\opt{A}{$\alpha>\beta$}
\opt{B}{$\beta>\gamma$}
\opt{C}{$\alpha<\gamma$}
\opt{D}{none of these}
\end{qoptions}
\qn{4} $\displaystyle\int_4^8\left(\sqrt{\sqrt{16x-64}+x}+\sqrt{x-\sqrt{16x-64}}\right)dx$
\begin{qoptions}
\opt{A}{$2\displaystyle\int_4^6\left(\sqrt{\sqrt{16x-64}+x}+\sqrt{x-\sqrt{16x-64}}\right)dx$}
\opt{B}{$2\displaystyle\int_8^{10}\left(\sqrt{\sqrt{16x-64}+x}+\sqrt{x-\sqrt{16x-64}}\right)dx$}
\opt{C}{$2\displaystyle\int_7^9\left(\sqrt{\sqrt{16x-64}+x}+\sqrt{x-\sqrt{16x-64}}\right)dx$}
\opt{D}{16}
\end{qoptions}
\qn{5} If M is the skew symmetric matrix of order $n\times n$, then
\begin{qoptions}
\opt{A}{$\det(M-I)=\det(M+I)$}
\opt{B}{$\det(M-I)=-\det(M+I)$ if n is odd}
\opt{C}{$\det(M-I)=\det(M+I)$ if n is even}
\opt{D}{$\det(I-M)(I+M)^{-1}=1$, $\det(I+M)\neq0$}
\end{qoptions}
\qn{6} If the unit vectors $\vec a$ and $\vec b$ are inclined at an angle $2\theta$ such that $|\vec a-\vec b|<1$ and $0\leqslant\theta\leqslant\pi$ then $\theta$ lies in the interval
\begin{qoptions}
\opt{A}{$[0,\pi/6)$}
\opt{B}{$(5\pi/6,\pi]$}
\opt{C}{$(\pi/6,\pi/2]$}
\opt{D}{$[\pi/2,5\pi/6)$}
\end{qoptions}
\sect{SECTION-I(ii)}
\qn{7} Match the List$-$I \& List$-$II\par
{\footnotesize\setlength{\tabcolsep}{1.5pt}\noindent\begin{tabular}{|c|L{4.4cm}|c|L{1.3cm}|} \hline & \multicolumn{1}{c|}{List-I} & & \multicolumn{1}{c|}{List-II}\\ \hline
(I) & A straight line with negative slope passing through (1, 4) meets the coordinate axes at A and B. The minimum length of OA + OB, O being the origin, is & (P) & $5\sqrt2$\\ \hline
(II) & If the point P is symmetric to the point Q(4, $-1$) with respect to the bisector of the first quadrant, then the length of PQ is & (Q) & $3\sqrt2$\\ \hline
(III) & On the portion of the straight line x + y = 2 between the axis a square is constructed away from the origin, with this portion as one of its sides. If d denotes the perpendicular distance of a side of this square from the origin then the maximum value of d is & (R) & 9/2\\ \hline
(IV) & If the parametric equation of a line is given x = 4 + $\lambda\sqrt2$ and y = $-1+\sqrt2\lambda$ where $\lambda$ is a parameter, then the intercept made by the line on the x-axis is & (S) & 9\\ \hline
\end{tabular}}\par
\begin{qoptions}
\opt{A}{I $\rightarrow$ R;II $\rightarrow$ P;III $\rightarrow$ Q;IV $\rightarrow$ S}
\opt{B}{I $\rightarrow$ S;II $\rightarrow$ P;III $\rightarrow$ Q;IV $\rightarrow$ R}
\opt{C}{I $\rightarrow$ S;II $\rightarrow$ R;III $\rightarrow$ Q;IV $\rightarrow$ P}
\opt{D}{I $\rightarrow$ P;II $\rightarrow$ S;III $\rightarrow$ Q;IV $\rightarrow$ R}
\end{qoptions}
\qn{8} Match the List$-$I \& List$-$II\par
{\footnotesize\setlength{\tabcolsep}{1.5pt}\noindent\begin{tabular}{|c|L{4.4cm}|c|L{1.0cm}|} \hline & \multicolumn{1}{c|}{List-I} & & \multicolumn{1}{c|}{List-II}\\ \hline
(I) & The smallest integer greater than $\dfrac{1}{\log_3\pi}+\dfrac{1}{\log_4\pi}$ is & (P) & 10\\ \hline
(II) & Let $3^a=4$, $4^b=5$, $5^c=6$, $6^d=7$, $7^e=8$, and $8^f=9$. Then the value of the product (abcdef) is & (Q) & 3\\ \hline
(III) & Characteristic of the logarithm of 2008 to the base 2 is & (R) & 1\\ \hline
(IV) & If $\log_2(\log_2(\log_3x))=\log_2(\log_3(\log_2y))=0$, then the value of $(x-y)$ is & (S) & 2\\ \hline
\end{tabular}}\par
\begin{qoptions}
\opt{A}{I $\rightarrow$ R;II $\rightarrow$ S;III $\rightarrow$ P;IV $\rightarrow$ Q}
\opt{B}{I $\rightarrow$ S;II $\rightarrow$ Q;III $\rightarrow$ P;IV $\rightarrow$ R}
\opt{C}{I $\rightarrow$ Q;II $\rightarrow$ S;III $\rightarrow$ P;IV $\rightarrow$ R}
\opt{D}{I $\rightarrow$ Q;II $\rightarrow$ R;III $\rightarrow$ P;IV $\rightarrow$ S}
\end{qoptions}
\qn{9} Match the following List-I with List-II\par
{\footnotesize\setlength{\tabcolsep}{1.5pt}\noindent\begin{tabular}{|c|L{4.6cm}|c|L{0.6cm}|} \hline & \multicolumn{1}{c|}{List-I} & & \multicolumn{1}{c|}{List-II}\\ \hline
(I) & $\displaystyle\lim_{x\to0}\dfrac{1-\cos2x}{e^{x^2}-e^x+x}$ equals to & (P) & 1\\ \hline
(II) & If the value of $\displaystyle\lim_{x\to0^+}\left(\dfrac{(3/x)+1}{(3/x)-1}\right)^{1/x}$ can be expressed in the form of $e^{p/q}$, where p and q are relative prime, then (p + q) is equal to & (Q) & 2\\ \hline
(III) & $\displaystyle\lim_{x\to0}\dfrac{\tan^3x-\tan^3}{x^5}$ equals to & (R) & 4\\ \hline
(IV) & \resizebox{4.5cm}{!}{$\displaystyle\lim_{x\to0}\frac{x+2\sin x}{\sqrt{x^2+2\sin x+1}-\sqrt{\sin^2x-x+1}}$} & (S) & 5\\ \hline
\end{tabular}}\par
\begin{qoptions}
\opt{A}{I $\rightarrow$ R;II $\rightarrow$ S;III $\rightarrow$ P;IV $\rightarrow$ Q}
\opt{B}{I $\rightarrow$ Q;II $\rightarrow$ S;III $\rightarrow$ S;IV $\rightarrow$ P}
\opt{C}{I $\rightarrow$ R;II $\rightarrow$ P;III $\rightarrow$ S;IV $\rightarrow$ Q}
\opt{D}{I $\rightarrow$ Q;II $\rightarrow$ S;III $\rightarrow$ P;IV $\rightarrow$ R}
\end{qoptions}
\qn{10} Let $\Pi$ be the plane parallel to y-axis and containing the points (1, 0, 1) and (3, 2, $-1$). A $\equiv$ (4, 0, 0) and B $\equiv$ (6, 0, $-2$) are two points and P $\equiv$ ($x_0,y_0,z_0$) is a variable point on the plane $\Pi=0$.\par
Match List$-$I with List$-$II and select the correct answer using the code given below the list.\par
{\footnotesize\setlength{\tabcolsep}{1.5pt}\noindent\begin{tabular}{|c|L{4.6cm}|c|L{0.7cm}|} \hline & \multicolumn{1}{c|}{List-I} & & \multicolumn{1}{c|}{List-II}\\ \hline
(I) & If the equation of the plane $\Pi=0$ is x + ay + bz = c, then $|a+b+c|$ is & (P) & 16\\ \hline
(II) & If (PA + PB) is minimum, then $|4x_0+y_0+2z_0|$ is & (Q) & 12\\ \hline
(III) & If $|PA+PB|\in\left[0,\sqrt N\right)$, then N is & (R) & 3\\ \hline
(IV) & If the reflection of the line AB in the plane $\Pi=0$ is $\dfrac{x-2}{1}=\dfrac{y-\alpha}{0}=\dfrac{z+\beta}{-1}$, then $(\alpha^4+\beta^4)$ & (S) & 8\\ \hline
 & & (T) & 10\\ \hline
\end{tabular}}\par
\begin{qoptions}
\opt{A}{I $\rightarrow$ R;II $\rightarrow$ Q;III $\rightarrow$ P;IV $\rightarrow$ S}
\opt{B}{I $\rightarrow$ Q;II $\rightarrow$ R;III $\rightarrow$ S;IV $\rightarrow$ P}
\opt{C}{I $\rightarrow$ R;II $\rightarrow$ Q;III $\rightarrow$ S;IV $\rightarrow$ P}
\opt{D}{I $\rightarrow$ R;II $\rightarrow$ S;III $\rightarrow$ Q;IV $\rightarrow$ P}
\end{qoptions}
\sect{SECTION-II}
\qn{1} In the coordinate plane, the region M consists of all the points (x, y) satisfying the inequalities $y\geq0$, $y\leq x\,y$ and $y\leq2-x$ simultaneously. The region N, which varies with the parameter t, consists of all the points (x, y) satisfying the inequalities $t\leq x\leq t+1$ and $0\leq t\leq1$ simultaneously. The area of M$\cap$N is $-k_1t^2+k_2t+\dfrac12$, then $k_1+k_2$ is \_\_\_\_\_\_ .
\qn{2} The slope of the tangent to the curve represented by $x=t^2+3t-8$ and $y=2t^2-2t-5$ at the point M(2, 1) is $\left(\dfrac\lambda7\right)$. then value of $\lambda$ is ?
\qn{3} Let $(1+3x+4x^2)=a_0+a_1x+a_2x^2+\ldots+a_{20}x^{20}$. Then value of $a_3$ is $-$
\qn{4} Find the value of $f\left(\dfrac\pi6\right)$, where $f(\theta)=\left|\begin{array}{ccc}\cos^2\theta&\cos\theta\sin\theta&-\sin\theta\\ \cos\theta\sin\theta&\sin^2\theta&\cos\theta\\ \sin\theta&-\cos\theta&0\end{array}\right|$
\qn{5} If $\alpha,\beta,\gamma$ are such that $\alpha+\beta+\gamma=4$, $\alpha^2+\beta^2+\gamma^2=6$, $\alpha^3+\beta^3+\gamma^3=8$, then the value of $[\alpha^4+\beta^4+\gamma^4]$ must be equal to (where [.] denotes the greatest integer function)
\qn{6} Find the Coefficient of $x^{103}$ in $(1+x+x^2+x^3+x^4)^{199}(x-1)^{201}$.
\qn{7} If $\left(1+\dfrac1x\right)^{x+1}=\left(1+\dfrac1{2021}\right)^{2021}$, then $|x|$ is equal to
\qn{8} If two events A and B are such that $P(\overline A)=0.3$, $P(B)=0.4$, $P(A\cap\overline B)=0.5$, then the value of $P(B/(A\cap\overline B))$ is
\raggedcolumns
\papersec{2}
\partsec{1}{PHYSICS}
\sect{SECTION-I (i)}
\qn{1} The refractive index of the medium within a certain region, x $>$ 0, y $>$ 0, changes with y. A thin light ray travelling in the x-direction strikes the medium at right angles and moves through the medium along a circular arc. How does the refractive index depend on y ?\par
[R : Radius of path \& $\eta_0$ : Initial refractive index]\par
\qfig{p2_phy_s1_q01_1}
\begin{qoptions}
\opt{A}{$\dfrac{2R}{R-y}\eta_0$}
\opt{B}{$\dfrac{R}{R-y}\dfrac{\eta_0}{2}$}
\opt{C}{$\dfrac{R+y}{R}\eta_0$}
\opt{D}{$\dfrac{R}{R-y}\eta_0$}
\end{qoptions}
\qn{2} A mass is suspended from a wire and is pulled along the length of wire, resulting in oscillations of time period $T_1$. The same mass is next attached to wire of the same material and length but double the cross-sectional area. The time period this time is $T_2$. Then $T_1/T_2$ is equal to
\begin{qoptions}
\opt{A}{1 : 2}
\opt{B}{2 : 1}
\opt{C}{1 : $\sqrt2$}
\opt{D}{$\sqrt2$ : 1}
\end{qoptions}
\qn{3} Initially, $C_1$ is charged but $C_2$ is uncharged. When switch is closed, charges on $C_1$ and $C_2$ are\par
\qfig{p2_phy_s1_q03_1}
\begin{qoptions}
\opt{A}{20 $\mu$C, 0}
\opt{B}{10 $\mu$C, 10 $\mu$C}
\opt{C}{30 $\mu$C, 10 $\mu$C}
\opt{D}{0, 20 $\mu$C}
\end{qoptions}
\qn{4} A standing wave is produced in a vapour of an element of atomic mass 254 (g mol$^{-1}$) at 400 K. Nodes are found to be 6.77 cm apart when frequency of source was 1000 Hz. The vapour used for experiment is (R = 8.31 JK$^{-1}$mol$^{-1}$)
\begin{qoptions}
\opt{A}{monatomic}
\opt{B}{diatomic}
\opt{C}{polyatomic}
\opt{D}{none of these}
\end{qoptions}
\sect{SECTION-I (ii)}
\qn{5} A particle of mass m is moving in a field where the potential energy is given by U(x) = $U_0$ (1 $-$ cos ax), where $U_0$ and a are constants and x is the displacement from mean position. Then (for small oscillation)
\begin{qoptions}
\opt{A}{the time period is T = $2\pi\sqrt{\dfrac{m}{aU_0}}$}
\opt{B}{the speed of particle is maximum at x = 0}
\opt{C}{the amplitude of oscillations is $\dfrac\pi a$}
\opt{D}{the time period is T = $2\pi\sqrt{\dfrac{m}{a^2U_0}}$}
\end{qoptions}
\qn{6} In an interference arrangement similar to Young's double slit experiment, the slits $S_1$ and $S_2$ are illuminated with coherent microwave sources, each of frequency $10^6$Hz. The sources are synchronized to have zero phase difference. The slits are separated by a distance d = 150.0 m. The intensity I($\theta$) is measured as a function of $\theta$, where $\theta$ is defined as shown. If $I_0$ is maximum intensity, then I($\theta$) for $0\leq\theta\leq90^\circ$ is given by\par
\qfig{p2_phy_s2_q06_1}
\begin{qoptions}
\opt{A}{I ($\theta$) = $\dfrac{I_0}{2}$ for $\theta=30^\circ$}
\opt{B}{I ($\theta$) = $\dfrac{I_0}{4}$ for $\theta=90^\circ$}
\opt{C}{I($\theta$) = $I_0$ for $\theta=0^\circ$}
\opt{D}{I($\theta$) is constant for all values of $\theta$}
\end{qoptions}
\qn{7} A uniform square plate of mass m and edge a initially at rest starts rotating about one of the edge under the action of a constant torque $\tau$. Then at the end of the 5$^{\text{th}}$ sec after start
\begin{qoptions}
\opt{A}{angular momentum is equal to 5$\tau$}
\opt{B}{kinetic energy is equal to $\dfrac{75\tau^2}{ma^2}$}
\opt{C}{angular momentum is equal to 2.5$\tau$}
\opt{D}{kinetic energy is equal to $\dfrac{75\tau^2}{2ma^2}$}
\end{qoptions}
\qn{8} A particle of mass m start moving at t = 0 due to a force F = $F_0\cos\omega t$, where $F_0$ and $\omega$ are constant. Particle will stop first time at t = $t_0$ and distance travel in this time interval is S, then
\begin{qoptions}
\opt{A}{$T_0=\dfrac\pi\omega$}
\opt{B}{$t_0=\dfrac{\pi}{2\omega}$}
\opt{C}{$S=\dfrac{2F_0}{m\omega^2}$}
\opt{D}{$S=\dfrac{F_0}{m\omega^2}$}
\end{qoptions}
\qn{9} An ideal gas under goes a thermodynamic cycle as shown in figure. Which of the following graphs represents the same cycle ?\par
\qfig{p2_phy_s2_q09_1}
\begin{qoptions}
\opt{A}{\qfigi{p2_phy_s2_q09_A}}
\opt{B}{\qfigi{p2_phy_s2_q09_B}}
\opt{C}{\qfigi{p2_phy_s2_q09_C}}
\opt{D}{\qfigi{p2_phy_s2_q09_D}}
\end{qoptions}
\qn{10} In the given circuit, initially the charge on capacitor is $Q_0$. At t = 0, the switch S is closed. Which of the following statement(s) is/are correct ?\par
\qfig{p2_phy_s2_q10_1}
\begin{qoptions}
\opt{A}{Maximum current through the inductor is $\dfrac{Q_0}{\sqrt{LC}}$.b}
\opt{B}{Charge on the capacitor is zero at t = $\dfrac\pi2\sqrt{LC}$.}
\opt{C}{Current through the inductor is maximum at t = $\dfrac\pi2\sqrt{LC}$.}
\opt{D}{Maximum magnetic energy can be stored in the inductor is $\dfrac{Q_0^2}{2C}$.}
\end{qoptions}
\sect{SECTION-III}
\qn{1} A vertical hollow cylinder contains an ideal gas. The gas is enclosed by a 5 kg movable piston with area of cross-section $5\times10^{-3}$ m$^2$. The gas is slowly heated from 300 K to 350 K and the piston rises by 0.1 m. The piston is now clamped at this position and the gas is cooled back to 300 K. The difference between the heat energy added during heating process and heat energy lost during cooling process is 11x. Find value of x. [Take $p_0=10^5$ N/m$^2$]
\qn{2} Graph 1 and 2 represent the spectral emissive power of a body at temperature $T_1$ and $T_1$. The ratio of area under curve 2 and curve 1 is.\par
\qfig{p2_phy_s3_q02_1}
\qn{3} In the diagram there are four conducting plates A, B, C and D placed parallel to each other at equal separation L. If plate C starts moving towards plate B with velocity v. Find the current (in mA) flowing in the wire connecting A and D. (assume all other plates to be fixed)\par
(Given: $q_2=2\mu$C, $q_3=3\mu$C, v = 30 m/s, L = 0.05 m)\par
\qfig{p2_phy_s3_q03_1}
\qn{4} Water is flowing in varying cross-section pipe. The areas of cross-sections 1, 2 and 3 are 1 cm$^2$, 2 cm$^2$ and A cm$^2$, respectively. Water levels are shown in different vertical tubes of manometers. The speed of water at cross-section 3 is $\dfrac1{\sqrt x}$ m/s. The value of x is\par
\qfig{p2_phy_s3_q04_1}
\qn{5} If the K$_\alpha$ radiation of Mo (Z = 42) has a wavelength of 0.71 \AA, calculate wavelength of the corresponding radiation of Cu. i.e. K$_\alpha$ for Cu (Z = 29). Find the value of nearest integer value (in \AA) of wavelength.
\qn{6} A particle is suspended by a light vertical inelastic string of length $\ell$ from a fixed support. At its equilibrium position, it is projected horizontally with a speed $\sqrt{6g\ell}$. What is the ratio of the tension in the string in its horizontal position to that in string, when the particle is vertically above the point of support?
\qn{7} The electric field on two sides of a charged plate is shown in the figure. If the charge density on the plate is n$\varepsilon_0$, then n is\par
\qfig{p2_phy_s3_q07_1}
\qn{8} A hydrogen like atom (atomic number z) is in a higher excited state of quantum number n. The excited state by successively emitting two photons of energy 10.2 eV and 17.0 eV, respectively. Alternative, the atom from the same excited by successively emitting two photons of energies 4.25 eV and 5.95 eV, respectively. Determine the value of n. (Take, ionisation energy of H-atom = 13.6 eV)
\raggedcolumns
\partsec{2}{CHEMISTRY}
\sect{SECTION-I (i)}
\qn{1} For [HgI$_3$]$^-$ ion geometry and magnetic behaviour is respectively:
\begin{qoptions}
\opt{A}{triangular bipyramidal, paramagnetic}
\opt{B}{triangular pyramidal, diamagnetic}
\opt{C}{triangular planar, paramagnetic}
\opt{D}{triangular planar, diamagnetic}
\end{qoptions}
\qn{2} \qfigi[1.0]{p2_chem_s1_q02_1}\par
The product (A) is
\begin{qoptions}
\opt{A}{\qfigi{p2_chem_s1_q02_A}}
\opt{B}{\qfigi{p2_chem_s1_q02_B}}
\opt{C}{\qfigi{p2_chem_s1_q02_C}}
\opt{D}{\qfigi{p2_chem_s1_q02_D}}
\end{qoptions}
\qn{3} For most reactions, rate of a reaction increases about two times when its temperature is raised by 10 K. This fact is due to
\begin{qoptions}
\opt{A}{increase in number of collisions per unit volume}
\opt{B}{increase in number of molecules having energy $E_A$ or more}
\opt{C}{lowering of activation energy}
\opt{D}{lowering of pre-exponential factor in Arrhenius equation}
\end{qoptions}
\qn{4} At constant volume, for a fixed amount of a gas the pressure of the gas increase with rise of temperature due to
\begin{qoptions}
\opt{A}{increase in average molecular speed}
\opt{B}{increased rate of collision amongst molecules}
\opt{C}{increase in molecular attraction}
\opt{D}{decrease in mean free path}
\end{qoptions}
\sect{SECTION-I (ii)}
\qn{5} Which of the following statements are correct ?
\begin{qoptions}
\opt{A}{Copper metal can reduce Fe$^{2+}$ in acidic medium.}
\opt{B}{The absolute value of standard electrode potential can be determined experimentally}
\opt{C}{In a Daniell cell, electrons flow from zinc electrode to copper electrode outside the cell}
\opt{D}{Cell voltage is an intensive quantity.}
\end{qoptions}
\qn{6} Which of the following show the correct order of nucleophilicity ?
\begin{qoptions}
\opt{A}{CH$_3$S$^-$ $>$ CH$_3$O$^-$}
\opt{B}{(CH$_3$)$_3$ N $>$ (CH$_3$)$_3$ $\ddot{\text{P}}$}
\opt{C}{CN$^-$ $>$ Cl$^-$}
\opt{D}{$\overset{-}{\text{O}}$C$_2$H$_5$ $>$ $\overset{-}{\text{O}}$H}
\end{qoptions}
\qn{7} Which of the following is/are more acidic than phenol ?
\begin{qoptions}
\opt{A}{\qfigi{p2_chem_s2_q07_A}}
\opt{B}{\qfigi{p2_chem_s2_q07_B}}
\opt{C}{\qfigi{p2_chem_s2_q07_C}}
\opt{D}{\qfigi{p2_chem_s2_q07_D}}
\end{qoptions}
\qn{8} Select molecules(s) in which all the $\pi$-orbitals are in same plane.
\begin{qoptions}
\opt{A}{C$_6$H$_6$}
\opt{B}{O = C = C = C = O}
\opt{C}{H$_2$C = C = CH$_2$}
\opt{D}{CH$_2$ = CH $-$ CH = CH$_2$}
\end{qoptions}
\qn{9} In which of the following Pauli rule is violated ?
\begin{qoptions}
\opt{A}{1s$^2$, 2s$^2$, 2p$^7$}
\opt{B}{1s$^1$, 2s$^1$, 2p$^6$}
\opt{C}{\qfigi[0.75]{p2_chem_s2_q09_C}}
\opt{D}{\qfigi[0.75]{p2_chem_s2_q09_D}}
\end{qoptions}
\qn{10} Synergic bonding in metal carbonyl complexes [M = Metal]:
\begin{qoptions}
\opt{A}{decreases C $-$ O bond strength}
\opt{B}{increases M $-$ C bond length}
\opt{C}{increases bond order for M $-$ C bond}
\opt{D}{increases paramagnetism of complex}
\end{qoptions}
\sect{SECTION-III}
\qn{1} Find the total number of electron in Cr which have (n $\times$ l) = zero
\qn{2} Total number of exchanges possible for 3d$^4$ electronic configuration are?
\qn{3} Total number of bonds in PCl$_5$ which are at 90$^\circ$ to each other, are
\qn{4} Consider octahedral complex K$_2$[Ni(EDTA)].\par
Total no. of 5 member rings in compound = X\par
Total number of geometrical isomers possible for compound = Y\par
Represent your answer as a total of X and Y.
\qn{5} Consider a titration of potassium dichromate solution with acidified Mohr's salt solution using diphenylamine as indicator. The amount of Mohr's salt required per mole of dichromate is
\qn{6} The number of atoms in a single face-centred cubic unit cell is
\qn{7} Total no. of geometrical isomers possible for 1, 2, 4-tribromocyclohexane are?
\qn{8} Total no. of stereoisomers possible for 1, 2-dimethylcyclopropane are?
\raggedcolumns
\partsec{3}{MATHEMATICS}
\sect{SECTION-I (i)}
\qn{1} If the area bounded by f (x) = $\sqrt{1+x^2}$ (x $>$ 0) the line y = x, y axis and y = $-$ x + a (a $>$1) is k, then area bounded by the graph of $f^{-1}$(x), the line y = x between the lines y = $-$ x + 1 and y = $-$ x + a is
\begin{qoptions}
\opt{A}{$\dfrac{k-1}{4}$}
\opt{B}{$\dfrac{2k-1}{2}$}
\opt{C}{k}
\opt{D}{$\dfrac{4k-1}{4}$}
\end{qoptions}
\qn{2} In a acute angle triangle ABC, the foot of the perpendicular from A divide the opposite side into two parts of length 3 and 17 units and $\tan A=\dfrac{22}{7}$, the $\left(\dfrac{\text{Area of }\Delta ABC}{22}\right)$ is
\begin{qoptions}
\opt{A}{$\dfrac52$}
\opt{B}{5}
\opt{C}{10}
\opt{D}{20}
\end{qoptions}
\qn{3} The digits of three digit number N are in A.P. If sum of the digits is 15 and the number obtained by reversing the digits of the number is 594 less the original number, then $\dfrac{1000}{N-252}$ is equal to :
\begin{qoptions}
\opt{A}{$\dfrac56$}
\opt{B}{$\dfrac53$}
\opt{C}{0.06}
\opt{D}{0.03}
\end{qoptions}
\qn{4} A flagstaff stands vertically on a pillar, the height of the flagstaff being double the height of the pillar. A man on the ground at a distance finds that both the pillar and the flagstaff subtend equal angles at his eyes. The ratio of the height of the pillar and the distance of the man from the pillar, is
\begin{qoptions}
\opt{A}{$\sqrt3:1$}
\opt{B}{1 : 3}
\opt{C}{$1:\sqrt3$}
\opt{D}{$\sqrt3:2$}
\end{qoptions}
\sect{SECTION-I (ii)}
\qn{5} Let $\vec a,\vec b,\vec c$ be non-coplanar vectors and $\vec p=\dfrac{\vec b\times\vec c}{[\vec a\,\vec b\,\vec c]}$, $\vec q=\dfrac{\vec c\times\vec a}{[\vec b\,\vec c\,\vec a]}$, $\vec r=\dfrac{\vec a\times\vec b}{[\vec c\,\vec a\,\vec b]}$ then
\begin{qoptions}
\opt{A}{$\vec p.\vec a=1$}
\opt{B}{$\vec p.\vec a+\vec q\cdot\vec b+\vec r.\vec c=3$}
\opt{C}{$\vec p\cdot\vec a+\vec q\cdot\vec b+\vec r.\vec c=0$}
\opt{D}{none of these}
\end{qoptions}
\qn{6} 7 points are given in a plane (as shown in adjacent figure), then using these 7 points $-$\par
\qfig{p2_math_s2_q06_1}
\begin{qoptions}
\opt{A}{Number of lines that can be formed is 10.}
\opt{B}{Number of triangles that can be formed is 24.}
\opt{C}{Number of quadrilateral that can be formed is 6.}
\opt{D}{Maximum number of circles that can be formed is 24, where each circle contains atleast three of 7 points}
\end{qoptions}
\qn{7} The equations of four circles are $(x\pm a)^2+(y\pm a)^2=a^2$. The radius of a circle touching all the four circles is
\begin{qoptions}
\opt{A}{$(\sqrt2-1)$ a}
\opt{B}{$2\sqrt2a$}
\opt{C}{$(\sqrt2+1)$ a}
\opt{D}{$(2+\sqrt2)$ a}
\end{qoptions}
\qn{8} If f(x) is a cubic polynomial such that f(0) + 2 = 0 and f$'''$(0) = 6 \& f (x) = 0 posses three positive integral roots then
\begin{qoptions}
\opt{A}{$f'(1)=0$}
\opt{B}{$2f'(0)+f''(0)=0$}
\opt{C}{$2f'(0)+f''(0)=2$}
\opt{D}{$f'(2)=0$}
\end{qoptions}
\qn{9} If x + [y] + \{z\} = 1.1\hfill [$\cdot$] is G.I.F. and \{$\cdot$\} is fractional part\par
[x] + \{y\} + z = 2.2\par
\{x\} + y + [z] = 3.3 then
\begin{qoptions}
\opt{A}{x + y + z = 3.3}
\opt{B}{y $-$ 2x = 1}
\opt{C}{2(z + 1) = 5y}
\opt{D}{\{x\} + \{y\} + \{z\} = 0.3}
\end{qoptions}
\qn{10} If the sides a, b \& c of a triangle ABC are in G.P. then $\angle B$ can never be
\begin{qoptions}
\opt{A}{$\dfrac\pi6$}
\opt{B}{$\dfrac{5\pi}{6}$}
\opt{C}{$\dfrac\pi4$}
\opt{D}{$\dfrac{2\pi}{3}$}
\end{qoptions}
\sect{SECTION-III}
\qn{1} If S = 3.$\sqrt9.\sqrt{\sqrt{27}.\sqrt{\sqrt{\sqrt{81}\ldots}}}\ldots\infty$, then sum of digits in S is
\qn{2} Let $f(x)=(1+x)\ell n(1+x)-x-\dfrac{x^2}{4}$, $\forall x>-1$ and $f'(x)>0\ \forall\in x(0,\alpha)$, then value of [$\alpha$] is (where [ ] denotes greatest integer function)
\qn{3} Let f(x) be a derivable function such that f(0) = 1, f(2) = 2, f(4) = 1 and f(6) = $-$ 4. The graph of y = f$'$ (x) is given below. Then find the total number of points of local maximum/local minimum of y = f (x) ln (0, 10)\par
\qfig{p2_math_s3_q03_1}
\qn{4} The projection length of a variable vector $x\hat i+y\hat j+z\hat k$ on the vector $\vec p=\hat i+2\hat j+3\hat k$ is 6. Let $\ell$ be the minimum projection length of the vector $x^2\hat i+y^2\hat j+z^2\hat k$ on the vector $\vec p$, then the value of $\sqrt[3]{\ell^2+15^2}$ is
\qn{5} Non-zero vectors $\vec a,\vec b,\vec c$ satisfy $\vec a.\vec b=0$, $\left(\vec b-\vec a\right).\left(\vec b+\vec c\right)=0$ and $2\left|\vec b+\vec c\right|=\left|\vec b-\vec a\right|$. If $\vec a=\mu\vec b+4\vec c$ then the value of $\mu$ is
\qn{6} 3 couples have to be seated around a circle. Let p be the probability that no couple is together then the value of 30 p is
\qn{7} If f(0) = 1, f(2) = 3, f$'$(2) = 5 then the value of the definite integral $\displaystyle\int_0^1xf''(2x)dx$ is
\qn{8} A special die is so constructed that the probabilities of throwing 1, 2, 3, 4, 5 and 6 are (1 $-$ k)/6, (1 + 2k)/6, (1 $-$ k)/6, (1 + k)/6, (1 $-$ 2k)/6 and (1 + k)/6 respectively. If two such dice are thrown and the probability of getting a sum equal to 9 lies in $\left[\dfrac19,\dfrac29\right]$. Then find the number of integral solutions of k.
\end{multicols}
\end{document}
